\section{The Artist of the Future}

According to \textbf{Boris Mouravieff}, there are four fields of study to understand man's interior and exterior worlds: \textbf{Philosophy}, \textbf{Religion}, \textbf{Science}, and \textbf{Art}. The history of Western civilization has been the dominance of one or another of those forms in succession. We are currently in the age of science, in which the scientist has displaced the philosopher of the ancient world and the theologian of the medieval. Art, however, is the synthesis of the preceding three forms. That synthesis is the opposite of a revolution, which is based on division. So reaction is not a looking back, but a going forward. Hermetic philosophy can elucidate these stages. But only a crisis provides the opportunities for the manifestation of latent forces, assuming they still exist.

The Triad of Providence, Will, and Destiny are the powers that govern the manifested universe (see \textbf{Rene Guenon}, \emph{The Great Triad}). The different ages focused on one of these powers. The synthesis will be knowledge of the whole, i.e., the tetrad.
\textbf{Fabre d'Olivet} explains:


\begin{quotex}
These three powers --- \textbf{Providence}; \textbf{Man}, or more precisely the human kingdom; and \textbf{Destiny} --- together make up the universal ternary. Nothing escapes their action. Everything in the universe is subordinate to them --- everything, that is, except for God Himself, who embraces all three in his unfathomable unity so as to form the tetrad of the ancients, that immense quaternary which is all in all and apart from which there is nothing.
\end{quotex}


\paragraph{The Age of Philosophy}


\begin{quotex}
Providential law is the law of the divine man. He lives a life primarily of the intellect, which is governed by that law.
\flright{\textsc{Fabre d'Olivet}}
\end{quotex}


It is the philosopher who lives a life of the intellect, that is, a life of contemplation of Being. In Being, there are the essences or, as Guenon calls them, ``possibilities''. From that perspective, there is no difference between possibilities of manifestation and possibilities of non-manifestation. The philosopher seeks the Wisdom to live the good life. For him, it is simply a question of knowledge. If ignorance of what is good is the problem, then knowledge is the solution. Since Being has no beginning, he regards the world as Eternal.

Nevertheless, this wisdom is available to the few, since the bulk of men neither seek it, nor understand it when it is explained to them. That is why Aristotle can claim that most men are born slaves. That does not mean they live as chattel slaves, but rather that they need to be directed.

That is the natural order of things, so the best state is the aristocratic, or rule of the best\footnote{\url{https://gornahoor.net/?p=7730}}. That requires, however, that men could recognize the best. That is not the case, since the best are recognized today only in athletic competitions and beauty. That is why they are so popular today. Hence, the degeneration of the castes is described by Plato as the fall from aristocracy to democracy in several stages.

\paragraph{The Age of Religion}


The next age is that of religion. Here the focus moves to the Will. Ignorance, or defect of the intelligence, is not the problem, but rather sin, or defect of the Will.

Thus, \textbf{Augustine} finds a meaning deeper than what the philosophers see. For him, slavery is the fruit of sin. As the Fathers often pointed out, the sinner has as many masters as he has vices. The doctrine of ``original sin'' just means that all men are born as slaves in this sense.

The philosopher believed a man could become just by overcoming ignorance. However, the theologian sees that justification can only come from something transcendent. Hence, man becomes free only through a second birth into this higher state.

Therefore, the just, or well-ordered, state is governed by the cross and the eagle.

\paragraph{The Age of Science}


As the awareness of Being and Will faded, attention fell on Destiny, which is nature under necessity. The claims of the philosopher and the theologian could not be verified in nature. The ordering principle of Religion, which had maintained civilization for centuries, lost its force as the Will weakened.

Religious restrictions came to be seen as forms of slavery rather than as the path to liberation from slavery. There is no need to repeat what's been said many times before, but Science did have great success in forming, or deforming, a civilization. However, the limitations of empiricism as an explanatory principle, are becoming quite apparent, even if only to a few.

There are two conflicting directions. On the one hand, science represents the final stage of the degeneration of the castes, not even anticipated by Guenon and Evola. Whereas the philosopher and theologian were aware of transcendent states, for the scientist, man is just another mammal. The consequences of this are seldom made explicit even by the educated and intelligent classes. Of course, it may very well be that current day rulers are nothing but mammals. Esoteric tradition teaches that there are degrees of men, dependent on their dominant chakra. Moreover, it also teaches that there are anthropoids with a human exterior, but without an intellectual soul. This is actually becoming accepted by mainstream thinking.

The other direction is the advancement of technology, which has made living easier in many ways. That opens the possibility for general prosperity, but the verdict is still out on that. However, it also makes possible mass manipulation and control. The will lead to a form of slavery unknown to previous eras.

\paragraph{The Age of the Artist}


\begin{quotex}
As far back as the beginnings of recorded history we find evidence that the Tradition taught the way to cross this moat by knowledge of oneself and by working on oneself.
\flright{\textsc{Boris Mouravieff}, \emph{Gnosis Book Two}}

The wise man will rule the stars
\flright{\textit{Rosicrucian maxim}}
\end{quotex}


The Artist of the future will be a philosopher, knight, saint, and scientist. He will integrate the intellectual, psychic and instinctive spheres into a whole. That integration is the True Will, so the Artist must needs be the conscious creator of his life. This is liberation from necessity, as Guenon explains:

\begin{quotex}
In uniting itself to Providence and consciously collaborating with it, the human Will can become a counter-balance to destiny and finally neutralise it... the Will hard won by faith this shows that it is related to Providence is capable of enslaving Necessity itself, controlling Nature and producing miracles. \flright{\textit{Pythagorean maxim}}
\end{quotex}


On the principle that Atman is Brahman, the knowledge of the whole, then, begins with self-knowledge. This knowledge leads to the Real I and True Will. Then,

\begin{quotex}
Esoteric Tradition teaches that any civilization is none other than a projection of the consciousness of the `I's of elite man onto the exterior world.
\end{quotex}


Hence, the Artist will need to have a creative imagination that can visualize a new civilization and bring it into manifestation.


\flrightit{Posted on 2016-02-23 by Cologero}

\begin{center}* * *\end{center}

\begin{footnotesize}
\begin{sffamily}
\texttt{Matthew Anderson on 2016-02-24 at 21:57 said:}

I like that term, ``intentional sociopath''. Here are the writings of a modern day intentional sociopath: \url{http://vongremmler.blogspot.com/}

Richard (the author of the blog) died in July, but he would have liked it too. It describes him well.

Formally, his articles are about ``psychology'', and he often refers to his ideas as ``theories''. But he uses that term to get past people's ego-defenses. His insights come neither from a scientific nor a theoretical perspective. His knowledge is from direct experience. For anyone wishing to develop their emotional center, this would be the perfect material for that.

\hfill

\texttt{Mark Citadel on 2016-03-08 at 09:34 said:}

Precise and on point as always:

``But only a crisis provides the opportunities for the manifestation of latent forces''

I become more certain of this with each passing day. Greatness cannot emerge from a world of decadence, a world of inverted caste. It must rise from the smoke of chaos, for in such an environment those who do not wake from their slumber perish, and those who do open their eyes find the organic will of God before them. Civilization rises like a phoenix.

\hfill

\texttt{Mark Citadel on 2016-03-10 at 17:13 said:}

By the way, Cologero, I just posted up a review of René Guénon's `The Crisis of the Modern World'. It's hard to express just how much I appreciated this book. A sage of Reaction!

\url{http://citadelfoundations.blogspot.com/2016/03/rene-guenons-crisis-of-modern-world.html}


\end{sffamily}
\end{footnotesize}

